\Definition{1.1}{
  Let $P$ be a real-valued set function defined on the collection of subsets of the
  sample space $S$. The set function $P:S\to[0,1]$ is called a \emph{probability
  measure} if it has the following properties:
  \begin{enumerate}[label=(\arabic*)]
    \item $0\le P(B)$, $B\subseteq S$.
    \item $P(S)=1$.
    \item If $B_i\cap B_j=\varnothing$ for $i,j=1,2,\dots$, $i\ne j$
      (pairwise disjoint), where $B_i\subseteq S$, then
      \[
        P\!\left(\bigcup_{i=1}^{\infty} B_i\right)
        = \sum_{i=1}^{\infty} P(B_i).
      \]
  \end{enumerate}
}

\Definition{1.2}{%
Let $B_1$ and $B_2$ be two events defined on a sample space $S$. 
The \emph{conditional probability of event $B_1$ given event $B_2$} (has occurred) is 
denoted as $P(B_1 \mid B_2)$ and defined as
\[
  P(B_1 \mid B_2) = \frac{P(B_1 \cap B_2)}{P(B_2)}.
\]
}

\Definition{1.3}{%
Let $B_1$ and $B_2$ be two events defined on a sample space $S$. 
Events $B_1$ and $B_2$ are said to be \emph{independent} if and only if
\[
  P(B_1 \cap B_2) = P(B_1) P(B_2).
\]
If the events $B_1$ and $B_2$ are not independent, they are said to be 
\emph{dependent}.
}

Therefore if $B_1$ and $B_2$ are independent, then $P(B_1 \mid B_2) = P(B_1) P(B_2)$

\Definition{1.4}{%
A \emph{random variable} $X$ is a real-valued function defined on the sample 
space $S$, 
\[
  X : S \to \mathbb{R} = (-\infty, \infty),
\]
where there is an associated probability measure $P$ defined on $S$. 
Let $A$ be the range of $X$, 
\[
  A = \{x \mid X(s) = x,\, s \in S\}.
\]
The range $A$ is known as the \emph{space of $X$} or \emph{state space of $X$}.
}
